\section{Experiments}
\label{sec:experiments}

\subsection{Setup and Natural-Length Comparison}
\label{sec:experimental_setup}
\label{sec:natural_length}

We use Qwen2.5-7B-Instruct as the teacher and Qwen2.5-1.5B-Instruct as the student, with 881 fixed GSM8K training problems~\citep{cobbe2021training}.
Students undergo completion-only LoRA SFT over three seeds and are evaluated by final-answer accuracy and generated tokens on 1,269 fixed test problems, using greedy decoding with a 512-token cap.
The primary layer-17, TopK-64 SAE identifies eight short-associated and eight long-associated features.
Appendix~\ref{app:details} provides sampling, SAE, and training details.

In the initial equal-example comparison, natural-short supervision achieves 66.67\% accuracy versus 64.20\% for natural-long supervision, below the independently evaluated base student (67.45\%).
The 2.47-point difference is not significant after Holm correction ($p=0.0852$).
Unequal supervision-token volumes also prevent attributing this ordering to redundant content.

\subsection{Teacher Intervention}
\label{sec:teacher_intervention}

Development screening selects enhancement of eight short-associated directions at relative perturbation magnitude $\rho=0.30$, using random controls matched in feature count and perturbation norm.
We fix the intervention before final student evaluation and generate four candidates per training problem under each condition.
Table~\ref{tab:teacher_generation} includes all outputs, including incorrect and capped responses.

\begin{table}[htbp]
\centering
\small
\begin{tabular}{lrrr}
\toprule
Generation condition & Correct (\%) & Mean tokens & Cap hits \\
\midrule
Unmodified & 98.69 & 250.27 & 12 \\
SAE short enhancement & 98.89 & 217.07 & 1 \\
Matched random & 97.70 & 293.02 & 38 \\
\bottomrule
\end{tabular}
\caption{Teacher generation on 881 problems, with 3,524 outputs per condition. Cap hits count responses reaching 512 tokens.}
\label{tab:teacher_generation}
\end{table}

SAE enhancement shortens outputs by 13.27\% relative to unmodified generation.
Question-level bootstrap intervals give a mean length difference of $-33.20$ tokens (95\% CI: $[-35.24,-31.22]$) and a correctness difference of $+0.20$ percentage points ($[-0.43,+0.79]$).
These are post-hoc descriptive results; reused feature-analysis cohorts and the absence of a non-inferiority test preclude a claim of established accuracy preservation.

\subsection{Downstream Distillation}
\label{sec:student_results}

For each generation condition, we deduplicate candidates and select each problem's shortest answer-correct, non-capped solution.
Intersecting support with the initial natural-short dataset leaves 878 shared problems.
Natural-short supervision uses a different pipeline---quantile-band selection from 16 candidates rather than shortest-correct selection from four---so this comparison does not isolate steering.
Four data conditions, two budgets, and three seeds yield 24 adapters.
Token matching repeats complete solutions; it does not equalize optimizer steps or per-question weights.

\begin{table}[htbp]
\centering
\small
\begin{tabular}{lrr}
\toprule
Supervision & Equal examples & Equal target tokens \\
\midrule
Unmodified & 69.92 & 70.42 \\
SAE short enhancement & 70.84 & 70.34 \\
Matched random & 68.87 & 68.87 \\
Natural short & 66.93 & 66.96 \\
\bottomrule
\end{tabular}
\caption{Mean student accuracy (\%) over three SFT seeds. Target-token budgets are approximately matched. The base student scores 66.51\% in this follow-up run.}
\label{tab:student_accuracy}
% Retain the legacy anchor so external references do not become undefined.
\label{fig:student_accuracy}
\end{table}

Table~\ref{tab:student_accuracy} summarizes student accuracy.
The SAE-minus-unmodified accuracy difference is $+0.92$ points under equal examples (95\% CI: $[-1.34,+3.23]$; Holm-adjusted $p=0.8432$) and $-0.08$ under equal tokens ($[-2.18,+1.97]$; $p=0.9654$).
Intervals use 10,000 crossed bootstrap replicates over seeds and paired evaluation problems; Holm correction covers six SAE-versus-control comparisons.
Comparisons with unmodified and random generation remain inconclusive.
SAE exceeds natural-short supervision under both budgets (adjusted $p=0.0060$ and $0.0110$), but the different sampling and selection pipelines confound attribution.

Student output length falls from 277.14 to 245.98 tokens under equal examples (11.24\%) and from 276.55 to 245.31 under equal tokens (11.30\%).
The follow-up base averages 239.64 tokens, so shortening is relative to unmodified distillation, not the pretrained student.
Its 66.51\% accuracy also differs from the initial run's 67.45\%; the cause remains unresolved.

\paragraph{Scope.}
\label{sec:experiment_scope}
The results support transfer of shorter-output behavior, not verified accuracy gains or end-to-end compute savings.
Evidence is limited to one teacher--student pair, GSM8K, one SAE seed, and reused analysis cohorts.
Concise prompting, dense short--long steering, ASC, textual compression, and length-matched controls remain untested~\citep{azizi2026activation,wu2025concise}; the value of SAE decomposition beyond simpler alternatives remains open.

\subsection{Baseline Comparison Design (Proposed)}
\label{sec:baseline_design}

The completed comparisons establish how the selected intervention behaves relative to its original controls.
They leave open whether similar supervision can be obtained with simpler generation or selection methods.
Table~\ref{tab:baseline_design} specifies additional comparisons; no unreported performance values are assigned to them.
The existing unmodified-generation student already uses shortest-correct selection from four candidates and must not be relabeled as random-correct.
Likewise, the historical natural-short condition cannot substitute for a comparison using the same candidate pool and selection rule.

\begin{table}[t]
\centering
\small
\setlength{\tabcolsep}{4pt}
\begin{tabularx}{\linewidth}{@{}>{\raggedright\arraybackslash}p{0.22\linewidth}>{\raggedright\arraybackslash}X>{\raggedright\arraybackslash}p{0.15\linewidth}@{}}
\toprule
Comparison & Question addressed & Evidence status \\
\midrule
Base student & Does the training pipeline improve on the initial student? & Reported \\
Unmodified + shortest-correct & Does steering add value with selection held fixed? & Reported \\
Matched random SAE & Is the chosen direction more useful than a matched perturbation? & One selected set reported \\
Natural short & How does the historical length-band pipeline compare? & Reported; selection differs \\
Random-correct, same pool & What does shortest-correct selection add to ordinary correct sampling? & Proposed \\
Concise prompt & Can explicit brevity instructions achieve comparable results? & Proposed \\
Dense short--long direction & Does an SAE decomposition add value over a direct activation contrast? & Proposed \\
Answer-format direction & Can answer formatting or an answer-transition correlate explain shortening? & Proposed \\
ASC & How does the method compare with learned activation compression? & Proposed \\
Text compression / DAP & How does internal control compare with rewriting teaching traces? & Proposed \\
\bottomrule
\end{tabularx}
\caption{Baseline coverage and planned additions. Reported rows refer only to the completed exploratory study; proposed rows have no completed results in this manuscript. Additional independent random feature sets are also proposed.}
\label{tab:baseline_design}
\end{table}

\paragraph{Selection, prompting, and internal controls.}
Random-correct and shortest-correct supervision would use the same unmodified pool, with random selection fixed before training.
A concise prompt would request a complete solution with brief explanations while keeping the required final-answer format fixed.
A dense baseline would average within-problem short--long activation differences with equal problem weighting at the same intervention layer.
An answer-format baseline would instead use paired texts with the same mathematical steps and different answer headings, controlling aligned positions when estimating the direction.
The latter is a targeted alternative explanation suggested by the feature-context evidence in Section~\ref{sec:answer_marker_evidence}, not an established mechanism.
Additional random SAE sets would be drawn independently before evaluation, matching feature count and perturbation norm and reporting all sampled controls.

\paragraph{Published compression methods.}
ASC would be implemented according to a specified published version, retaining its CES objective and KL constraint~\citep{azizi2026activation}.
A plain difference-of-means direction would remain a separate baseline.
Text compression would rewrite the same teacher solutions before student SFT; a full difficulty-aware pipeline would be distinguished from a generic shortening prompt~\citep{wu2025concise}.
Any additional reference model, rewriting call, or data requirement would be disclosed and included in cost accounting.
These methods would be evaluated with the same student rather than compared through their papers' original benchmark scores.

\paragraph{Two complementary fairness questions.}
A fixed-candidate comparison would equalize the number of sampling opportunities, verification rules, and shortest-correct selection across generation methods.
We would report coverage and failures before constructing a common-support student comparison, since intersections can hide problems on which a method fails.
A separate fixed-generation-token comparison would assess usable-correct-trace yield under a common acquisition budget.
These budgets answer different questions and would not be presented as simultaneously equal.
Student experiments would retain separate equal-example and approximately equal-target-token analyses, documenting repeat counts and optimizer steps.

\paragraph{Length matching and calibration.}
Steering strengths and prompt variants would be calibrated on development problems to obtain comparable length ranges, with all tested settings disclosed.
Accuracy--length curves would accompany comparisons at a chosen compression level; equal hidden-state norms alone do not imply equal behavioral effects.
Within-problem length matching could provide a complementary conditional comparison, with its tolerance, residual imbalance, and retained support reported.
Because length is itself affected by the method, post-hoc matching would not alone establish a causal content effect.
Selection would use separate GSM8K training problems, and final student evaluation would retain the fixed cohort.
Repeated use of that cohort would remain explicit rather than being described as a newly untouched confirmation set.

\paragraph{Endpoints and interpretation.}
Primary downstream comparisons would report paired student accuracy with uncertainty, alongside student output length and actual training tokens.
Teacher endpoints would include all-output correctness, cap hits, coverage, and usable correct traces per generated million tokens.
End-to-end cost would include activation collection, SAE or direction fitting, candidate generation, verification, rewriting, and online intervention overhead.
Both total cost and marginal cost with a reusable dictionary would be reported.
An inconclusive accuracy difference would not establish non-inferiority, even if the method yields shorter outputs.
